\section{Conclusion}
\label{sec:conclusion}

We have presented a multi-operator knowledge graph pipeline that generates
cross-domain hypothesis candidates by composing domain-specific knowledge graphs
via bridge alignment.
Across three independent bio-chem domain pairs on real Wikidata-derived KGs,
we establish three findings:

\begin{enumerate}[noitemsep]
\item Bridge alignment is \textit{necessary and sufficient} for cross-domain path
  reachability; the effect amplifies non-linearly with KG size (42$\times$ more
  unique paths for a 9$\times$ node increase).

\item A relation-type drift filter eliminates semantic drift (25\%$\to$0\% drift rate)
  while preserving all qualitatively promising candidates, and generalises across
  independent domain pairs without parameter retuning.

\item Bridge entity structural diversity---not bridge count---is the primary driver
  of candidate yield variation (0.7$\times$ to 8.0$\times$ per bridge pair across subsets).
\end{enumerate}

The post-hoc relation semantics audit (Run~014) identifies a unified root cause
for all observed candidate weaknesses: the KG encodes relation types at the
class level, which the current drift filter cannot resolve without substrate-specificity
metadata.
This finding is the primary constraint on interpreting any candidate from the current
pipeline as a mechanistic hypothesis.

\paragraph{Future work priorities.}

\begin{enumerate}[noitemsep]
\item \textbf{Scale to production KGs}: test the pipeline on PrimeKG, Hetionet,
  or ROBOKOP to establish whether the alignment-dependent reachability
  mechanism holds at millions-of-node scale.

\item \textbf{Substrate-specificity metadata}: add a \texttt{produces\_class}
  distinction to \texttt{produces} edges (substrate-specific vs.\ class-level)
  and compartment annotations to nodes, enabling a next-generation filter that
  can detect class-level chemistry failures.

\item \textbf{Expert validation of Subsets B and C candidates}: the 39 (Subset~B)
  and 33 (Subset~C) filter-passing deep candidates are currently labelled
  ``promising'' by filter inference only.
  Independent domain-expert labeling would convert these from inferred to
  verified promising candidates.

\item \textbf{Depth-promotion}: develop a ranking mechanism that allows deep
  cross-domain candidates to surface in top-$k$ without sacrificing the scoring
  quality of well-supported 2-hop chains.

\item \textbf{Formal bridge dispersion metric}: quantify bridge dispersion as
  a measurable KG property (e.g., harmonic mean of bridge-entity degree in
  both domain subgraphs) to enable predictive KG construction guidelines.
\end{enumerate}
